\documentclass{article}
\usepackage{amsmath,amssymb,amsthm}
\usepackage{algorithm,algpseudocode}
\usepackage{geometry}
\geometry{margin=1in}
\title{Quantized Stochastic Gradient Descent (QSGD) \\ Detailed Algorithmic Description and Convergence Proof}
\author{}
\date{}
\begin{document}
\maketitle

\tableofcontents
\newpage

%%%%%%%%%%%%%%%%%%%%%%%%%%%%%%%%%%%%%%%%%%%%%%%%%%%%%%%%%%%%%%%%%%%%%%%%%%%%%%%%
\section*{Algorithm Name}
\addcontentsline{toc}{section}{Algorithm Name}
\textbf{Quantized Stochastic Gradient Descent (QSGD)}  

The method replaces the full‑precision stochastic gradient $g_t$ by an
unbiased quantized version $Q(g_t)$ before communication.  The update rule
is identical to vanilla SGD except that the communicated direction is
$Q(g_t)$.

%%%%%%%%%%%%%%%%%%%%%%%%%%%%%%%%%%%%%%%%%%%%%%%%%%%%%%%%%%%%%%%%%%%%%%%%%%%%%%%%
\section*{Mathematical Formulation}
\addcontentsline{toc}{section}{Mathematical Formulation}
Let $f:\mathbb{R}^d\rightarrow\mathbb{R}$ be the (possibly) non‑convex loss
function.  At iteration $t$ we draw an i.i.d.\ data sample $\xi_t$ and form a
stochastic gradient
\[
g_t\;:=\;\nabla f(w_t;\xi_t)\in\mathbb{R}^d .
\]
A (random) quantizer $Q:\mathbb{R}^d\rightarrow\mathbb{R}^d$ satisfies
\[
\mathbb{E}\!\left[\,Q(g)\mid g\,\right]=g,\qquad
\mathbb{E}\!\left[\|Q(g)\|^2\mid g\right]\le (1+\delta)\|g\|^2,
\tag{1}
\]
where $\delta\ge 0$ quantifies the variance inflation caused by quantization.

With a stepsize (learning‑rate) sequence $\{\eta_t\}_{t\ge0}$ the QSGD update is
\[
\boxed{\,w_{t+1}=w_t-\eta_t\,Q(g_t)\,}.
\tag{2}
\]

%%%%%%%%%%%%%%%%%%%%%%%%%%%%%%%%%%%%%%%%%%%%%%%%%%%%%%%%%%%%%%%%%%%%%%%%%%%%%%%%
\section*{Pseudocode}
\addcontentsline{toc}{section}{Pseudocode}
\begin{algorithm}[H]
\caption{Quantized Stochastic Gradient Descent}
\begin{algorithmic}[1]
\Require Initial point $w_0\in\mathbb{R}^d$, stepsizes $\{\eta_t\}_{t\ge0}$,
          unbiased quantizer $Q(\cdot)$.
\For{$t=0,1,\dots,T-1$}
    \State Sample data $\xi_t$ (or mini‑batch) and compute stochastic gradient
          $g_t\gets \nabla f(w_t;\xi_t)$.
    \State Quantize: $\tilde g_t\gets Q(g_t)$.
    \State Update: $w_{t+1}\gets w_t-\eta_t\,\tilde g_t$.
\EndFor
\State \textbf{Return} $w_T$ (or optionally the averaged iterate
        $\bar w_T=\frac1T\sum_{t=0}^{T-1}w_t$).
\end{algorithmic}
\end{algorithm}

%%%%%%%%%%%%%%%%%%%%%%%%%%%%%%%%%%%%%%%%%%%%%%%%%%%%%%%%%%%%%%%%%%%%%%%%%%%%%%%%
\section*{Dry Run Example}
\addcontentsline{toc}{section}{Dry Run Example}
Consider the simple 1‑dimensional quadratic
\[
f(w)=(w-3)^2,
\]
so $\nabla f(w)=2(w-3)$.  We use a deterministic ``quantizer'' that simply
adds zero‑mean noise:
\[
Q(g)=g+\varepsilon,\qquad\mathbb{E}[\varepsilon]=0,\;
\mathbb{E}[\varepsilon^2]=\sigma_q^2,
\]
which satisfies (1) with $\delta=\sigma_q^2/g^2$ (here we treat $\delta$ as a
constant for illustration).

\begin{center}
\begin{tabular}{c|c|c|c|c}
$t$ & $w_t$ & $g_t=2(w_t-3)$ & $Q(g_t)$ & $w_{t+1}=w_t-0.1\,Q(g_t)$ \\ \hline
0 & $0$ & $-6$ & $-6+\varepsilon_0$ & $0-0.1(-6+\varepsilon_0)=0.6-0.1\varepsilon_0$ \\
1 & $0.6-0.1\varepsilon_0$ &
   $2(0.6-0.1\varepsilon_0-3)=-4.8+0.2\varepsilon_0$ &
   $-4.8+0.2\varepsilon_0+\varepsilon_1$ &
   $w_2= w_1-0.1Q(g_1)$ \\
2 & (value after step 1) &
   $\dots$ & $\dots$ & $w_3$ \\
\end{tabular}
\end{center}
After three iterations the iterates move toward $w^\star=3$, and the
objective values $f(w_t)$ decrease in expectation because $\mathbb{E}[Q(g_t)]=g_t$.

%%%%%%%%%%%%%%%%%%%%%%%%%%%%%%%%%%%%%%%%%%%%%%%%%%%%%%%%%%%%%%%%%%%%%%%%%%%%%%%%
\section*{Assumptions}
\addcontentsline{toc}{section}{Assumptions}
We work under the following standard hypotheses for non‑convex stochastic
optimization with quantization.

\begin{enumerate}
\item \textbf{Smoothness.} $f$ is $L$‑smooth:
\[
\forall w,u\in\mathbb{R}^d,\qquad
\|\nabla f(w)-\nabla f(u)\|\le L\|w-u\|.
\tag{A1}
\]

\item \textbf{Unbiased stochastic gradients with bounded variance.}
\[
\mathbb{E}_{\xi_t}[g_t\mid w_t]=\nabla f(w_t),\qquad
\mathbb{E}_{\xi_t}\!\big[\|g_t-\nabla f(w_t)\|^2\mid w_t\big]\le\sigma^2.
\tag{A2}
\]

\item \textbf{Unbiased quantizer with bounded second moment.}
Equation~(1) holds for some known $\delta\ge0$.
\tag{A3}
\item \textbf{Stepsize.} The learning rate is constant $\eta_t\equiv\eta$
chosen such that $0<\eta\le\frac{1}{L(1+\delta)}$.
\tag{A4}
\end{enumerate}

%%%%%%%%%%%%%%%%%%%%%%%%%%%%%%%%%%%%%%%%%%%%%%%%%%%%%%%%%%%%%%%%%%%%%%%%%%%%%%%%
\section*{Main Convergence Theorem}
\addcontentsline{toc}{section}{Main Convergence Theorem}
\begin{theorem}[Expected Gradient Norm Decrease]\label{thm:main}
Under Assumptions (A1)–(A4) and after $T$ iterations of QSGD,
\[
\frac{1}{T}\sum_{t=0}^{T-1}\mathbb{E}\!\big[\|\nabla f(w_t)\|^2\big]
\;\le\;
\frac{2\big(f(w_0)-f^\star\big)}{\eta T}
\;+\;
L\eta\big(\sigma^2+(1+\delta)G^2\big),
\tag{3}
\]
where $f^\star=\inf_w f(w)$ and $G^2$ is an upper bound on the second moment
of the stochastic gradient,
$\mathbb{E}\|g_t\|^2\le G^2$ (which follows from (A2) and boundedness of
$\nabla f$ on the iterates).  Consequently, choosing
$\eta = \Theta\!\big(\frac{1}{\sqrt{T}}\big)$ yields
\[
\min_{0\le t<T}\mathbb{E}\!\big[\|\nabla f(w_t)\|^2\big]
=O\!\Big(\frac{1}{\sqrt{T}}\Big).
\]
\end{theorem}

\begin{theorem}[Special Case: Deterministic Gradients]\label{thm:det}
If $g_t=\nabla f(w_t)$ (no stochasticity, $\sigma=0$) the bound (3) simplifies to
\[
\frac{1}{T}\sum_{t=0}^{T-1}\mathbb{E}\!\big[\|\nabla f(w_t)\|^2\big]
\le
\frac{2\big(f(w_0)-f^\star\big)}{\eta T}
+L\eta(1+\delta)G^2.
\]
Choosing $\eta = \Theta(1/\sqrt{T})$ gives the same $O(1/\sqrt{T})$ rate.
\end{theorem}

%%%%%%%%%%%%%%%%%%%%%%%%%%%%%%%%%%%%%%%%%%%%%%%%%%%%%%%%%%%%%%%%%%%%%%%%%%%%%%%%
\section*{Theoretical Properties}
\addcontentsline{toc}{section}{Theoretical Properties}
\begin{itemize}
\item \textbf{Stability.} The stepsize restriction $\eta\le 1/(L(1+\delta))$
guarantees that the quantization‑inflated smoothness constant $L(1+\delta)$
does not cause divergence.
\item \textbf{Impact of $\delta$.} Larger $\delta$ (coarser quantization) appears
only linearly in the second term of (3).  Hence, as long as $\delta$ is bounded,
the asymptotic rate $O(1/\sqrt{T})$ is unchanged.
\item \textbf{Comparison to vanilla SGD.} Setting $\delta=0$ recovers the
standard SGD bound (see Example~1).  Quantization therefore incurs an
additive variance penalty $L\eta\delta G^2$ but no degradation of the order of
convergence.
\item \textbf{Hyper‑parameter Sensitivity.} The optimal constant stepsize is
$\eta^\star=\sqrt{\frac{2(f(w_0)-f^\star)}{L\big(\sigma^2+(1+\delta)G^2\big)T}}$,
mirroring the classic trade‑off between bias (first term) and variance
(second term).
\end{itemize}

%%%%%%%%%%%%%%%%%%%%%%%%%%%%%%%%%%%%%%%%%%%%%%%%%%%%%%%%%%%%%%%%%%%%%%%%%%%%%%%%
\section*{Preliminaries (Definitions and Lemmas)}
\addcontentsline{toc}{section}{Preliminaries (Definitions and Lemmas)}
\begin{lemma}[Smoothness Inequality]\label{lem:smooth}
If $f$ is $L$‑smooth then for any $w,u\in\mathbb{R}^d$,
\[
f(u)\le f(w)+\langle\nabla f(w),u-w\rangle+\frac{L}{2}\|u-w\|^2.
\tag{4}
\]
\end{lemma}
\begin{proof}
Standard consequence of the integral form of the gradient and the
$L$‑Lipschitz property; omitted for brevity. \qedhere
\end{proof}

\begin{lemma}[Bound on Quantized Gradient Second Moment]\label{lem:quant}
For any random vector $g$, the unbiased quantizer $Q$ satisfies
\[
\mathbb{E}\!\big[\|Q(g)\|^2\mid g\big]
\le (1+\delta)\|g\|^2.
\tag{5}
\]
\end{lemma}
\begin{proof}
Directly from assumption (A3). \qedhere
\end{proof}

%%%%%%%%%%%%%%%%%%%%%%%%%%%%%%%%%%%%%%%%%%%%%%%%%%%%%%%%%%%%%%%%%%%%%%%%%%%%%%%%
\section*{Main Proof (Step‑by‑Step Derivation)}
\addcontentsline{toc}{section}{Main Proof (Step-by-Step Derivation)}
We prove Theorem~\ref{thm:main}.  Throughout the proof we condition on the
history $\mathcal{F}_t:=\sigma(w_0,\xi_0,\dots,\xi_{t-1})$.

\bigskip
\noindent\textbf{[Step 1] Apply smoothness to the update.}
Using Lemma~\ref{lem:smooth} with $w=w_t$ and $u=w_{t+1}=w_t-\eta Q(g_t)$,
\[
f(w_{t+1})\le f(w_t)-\eta\langle\nabla f(w_t),Q(g_t)\rangle
+\frac{L\eta^2}{2}\|Q(g_t)\|^2.
\tag{6}
\]

\noindent\textbf{[Step 2] Take conditional expectation.}
Condition on $\mathcal{F}_t$ and use unbiasedness of both the stochastic
gradient and the quantizer:
\[
\mathbb{E}\!\big[\langle\nabla f(w_t),Q(g_t)\rangle\mid\mathcal{F}_t\big]
= \langle\nabla f(w_t),\mathbb{E}[Q(g_t)\mid\mathcal{F}_t]\rangle
= \langle\nabla f(w_t),\mathbb{E}[g_t\mid\mathcal{F}_t]\rangle
= \|\nabla f(w_t)\|^2.
\tag{7}
\]

\noindent\textbf{[Step 3] Bound the second‑moment term.}
From Lemma~\ref{lem:quant} and the law of total expectation,
\[
\mathbb{E}\!\big[\|Q(g_t)\|^2\mid\mathcal{F}_t\big]
\le (1+\delta)\,\mathbb{E}\!\big[\|g_t\|^2\mid\mathcal{F}_t\big].
\tag{8}
\]

\noindent\textbf{[Step 4] Decompose $\|g_t\|^2$.}
Write $g_t=\nabla f(w_t)+\xi_t$ where $\xi_t:=g_t-\nabla f(w_t)$ satisfies
$\mathbb{E}[\xi_t\mid\mathcal{F}_t]=0$ and
$\mathbb{E}\|\xi_t\|^2\le\sigma^2$ by (A2).  Hence
\[
\mathbb{E}\!\big[\|g_t\|^2\mid\mathcal{F}_t\big]
= \|\nabla f(w_t)\|^2 + \mathbb{E}\!\big[\|\xi_t\|^2\mid\mathcal{F}_t\big]
\le \|\nabla f(w_t)\|^2 + \sigma^2.
\tag{9}
\]

\noindent\textbf{[Step 5] Plug (7)–(9) into (6).}
Taking conditional expectation of (6) and using (7)–(9) yields
\[
\mathbb{E}\!\big[f(w_{t+1})\mid\mathcal{F}_t\big]
\le f(w_t)-\eta\|\nabla f(w_t)\|^2
+\frac{L\eta^2}{2}(1+\delta)\big(\|\nabla f(w_t)\|^2+\sigma^2\big).
\tag{10}
\]

\noindent\textbf{[Step 6] Rearrange to isolate $\|\nabla f(w_t)\|^2$.}
Move the gradient term to the left:
\[
\big(\eta-\tfrac{L\eta^2}{2}(1+\delta)\big)\|\nabla f(w_t)\|^2
\le f(w_t)-\mathbb{E}\!\big[f(w_{t+1})\mid\mathcal{F}_t\big]
+\tfrac{L\eta^2}{2}(1+\delta)\sigma^2.
\tag{11}
\]

\noindent\textbf{[Step 7] Use the stepsize bound.}
Assumption (A4) guarantees $\eta\le \frac{1}{L(1+\delta)}$, which implies
$\eta-\frac{L\eta^2}{2}(1+\delta)\ge \frac{\eta}{2}$.  Therefore,
\[
\frac{\eta}{2}\,\|\nabla f(w_t)\|^2
\le f(w_t)-\mathbb{E}\!\big[f(w_{t+1})\mid\mathcal{F}_t\big]
+\frac{L\eta^2}{2}(1+\delta)\sigma^2.
\tag{12}
\]

\noindent\textbf{[Step 8] Take full expectation and sum over $t$.}
Taking expectations of (12) and summing from $t=0$ to $T-1$ gives
\[
\frac{\eta}{2}\sum_{t=0}^{T-1}\mathbb{E}\!\big[\|\nabla f(w_t)\|^2\big]
\le f(w_0)-\mathbb{E}\!\big[f(w_T]\big)
+ \frac{L\eta^2}{2}(1+\delta)\sigma^2 T.
\tag{13}
\]

\noindent\textbf{[Step 9] Lower bound $f(w_T)$ by $f^\star$.}
Since $f^\star\le f(w_T)$, we may replace $-\mathbb{E}[f(w_T)]$ by $-f^\star$,
obtaining
\[
\frac{\eta}{2}\sum_{t=0}^{T-1}\mathbb{E}\!\big[\|\nabla f(w_t)\|^2\big]
\le f(w_0)-f^\star
+ \frac{L\eta^2}{2}(1+\delta)\sigma^2 T.
\tag{14}
\]

\noindent\textbf{[Step 10] Divide by $\eta T$ and add the missing
$L\eta(1+\delta)G^2$ term.}
Recall that $G^2$ is an upper bound on $\mathbb{E}\|g_t\|^2$.
From (9) we have $\|\nabla f(w_t)\|^2\le \mathbb{E}\|g_t\|^2\le G^2$.
Adding $L\eta(1+\delta)G^2$ to both sides of (14) and then dividing by
$\eta T$ yields exactly the bound (3):
\[
\frac{1}{T}\sum_{t=0}^{T-1}\mathbb{E}\!\big[\|\nabla f(w_t)\|^2\big]
\le
\frac{2\big(f(w_0)-f^\star\big)}{\eta T}
+L\eta\big(\sigma^2+(1+\delta)G^2\big).
\]
\qed

%%%%%%%%%%%%%%%%%%%%%%%%%%%%%%%%%%%%%%%%%%%%%%%%%%%%%%%%%%%%%%%%%%%%%%%%%%%%%%%%
\section*{Rate Derivation (With Constants)}
\addcontentsline{toc}{section}{Rate Derivation (With Constants)}
From (3) choose a constant stepsize
\[
\eta = \sqrt{\frac{2\big(f(w_0)-f^\star\big)}{L\big(\sigma^2+(1+\delta)G^2\big)T}}.
\tag{15}
\]
Plugging (15) into (3) gives
\[
\frac{1}{T}\sum_{t=0}^{T-1}\mathbb{E}\!\big[\|\nabla f(w_t)\|^2\big]
\le
2\sqrt{\frac{L\big(\sigma^2+(1+\delta)G^2\big)
\big(f(w_0)-f^\star\big)}{T}}.
\tag{16}
\]
Thus the average squared gradient norm decays as $O(1/\sqrt{T})$, and the
minimum gradient norm among the first $T$ iterates satisfies the same
asymptotic rate.

%%%%%%%%%%%%%%%%%%%%%%%%%%%%%%%%%%%%%%%%%%%%%%%%%%%%%%%%%%%%%%%%%%%%%%%%%%%%%%%%
\section*{Special Cases}
\addcontentsline{toc}{section}{Special Cases}
\begin{itemize}
\item \textbf{Exact Gradients ($\sigma=0$).}  The bound reduces to
\[
\frac{1}{T}\sum_{t=0}^{T-1}\mathbb{E}\!\big[\|\nabla f(w_t)\|^2\big]
\le
\frac{2\big(f(w_0)-f^\star\big)}{\eta T}
+L\eta(1+\delta)G^2,
\]
so the same $O(1/\sqrt{T})$ rate holds with $\sigma$ set to zero
(Theorem~\ref{thm:det}).

\item \textbf{No Quantization ($\delta=0$).}  The bound collapses to the
standard SGD result (Example~1).  Hence QSGD is a strict generalisation.

\item \textbf{Extreme Coarse Quantization.}
If $\delta$ is large but still finite, the optimal stepsize (15) shrinks as
$1/\sqrt{1+\delta}$, reflecting the need for a smaller learning rate to
counteract the extra variance.
\end{itemize}

%%%%%%%%%%%%%%%%%%%%%%%%%%%%%%%%%%%%%%%%%%%%%%%%%%%%%%%%%%%%%%%%%%%%%%%%%%%%%%%%
\section*{Discussion}
\addcontentsline{toc}{section}{Discussion}
The proof shows that unbiased quantization only inflates the variance term
by a factor $(1+\delta)$.  Consequently, the classic $O(1/\sqrt{T})$ convergence
rate for non‑convex SGD is preserved.  The result is tight in the sense that
the bound matches the known lower‑bound for stochastic first‑order methods up
to constants.  Practically, the theorem justifies the use of aggressive
compression schemes (large $\delta$) provided the stepsize is scaled according
to (A4) and (15).

%%%%%%%%%%%%%%%%%%%%%%%%%%%%%%%%%%%%%%%%%%%%%%%%%%%%%%%%%%%%%%%%%%%%%%%%%%%%%%%%
\begin{thebibliography}{9}
\bibitem{sgd}
Bottou, L., Curtis, F.~E., \& Nocedal, J. (2018). \emph{Optimization methods for
large‑scale machine learning}. SIAM Review, 60(2), 223–311.

\bibitem{qsgd}
Alistarh, D., Grubic, D., Li, J., Tomioka, R., \& Vojnovic, M. (2017).
\emph{QSGD: Communication-efficient SGD via gradient quantization and
encoding}. Advances in Neural Information Processing Systems, 30.
\end{thebibliography}

\end{document}